ここでは高被引用論文掲載誌の変化を高被引用でないグループのそれと比較する。G1-G4の当該論文の掲載数上位の雑誌一覧を表\ref{tab:G1topj}-\ref{tab:G4topj}に示す。各誌のインパクトファクターとその発表年についてはJournal Citation Reports  (JCR）[JCR1]から求めた。JCRではインパクトファクターの公開が1997からであるため、1946-2000年の累積世代以降についてのみ、累積世代の最終年のインパクトファクターを求めた。ここから、高被引用論文を持つ雑誌は必ずしもインパクトファクターが高い雑誌ではないことがわかる。一方、雑誌の構成においてはグループ間に特別な差異を認めらない。そこで、各グループでどの程度の雑誌の嗜好があるかを確認するため、全てのグループを合成した集団と各グループとの間でKullback-Leibler divergence（KLD）を求めたものを図14に示す。図14からは、高被引用論文は他のグループに比べ雑誌の嗜好が強いことが明らかになった。
